\section{An original-integral upper bound for the quadratic seven factor at two}
\label{sec:quadratic-seven-coefficient-two-upper-bound}

\paragraph{The actual family and its original coefficient.}
Put $K=\mathbb Q(\sqrt{-3})$, $\xi=\chi_{-7}$ and
$\epsilon_K=\chi_{-3}$.  Let $\delta$ be a Dirichlet character
of two-power conductor, including the trivial character.
Let $E/\mathbb Q_2$ be any finite extension containing its values,
with its original ring of integers $O$, and put
\[
 a=\delta\xi c_2^{-e},\qquad
 \mathscr D=(E/O)(a),\qquad b=[E:\mathbb Q_2],\qquad e\geq1.
 \tag{G189.1}
\]
Write $\mathcal S(K,\mathscr D)$ for full local zero at the unique
place above two and zero on every other finite inertia group.
No extra continuous character is included in (G189.1).
For $f\geq21$ and $e\geq1$ define the positive real number
\[
 \mathcal B(e,f)=
 \log\!\left(48\cdot10^e(e+1)(e+2)!e!\,
                              f^{e+3/2}\right).
 \tag{G189.2}
\]

\begin{theorem}[AC-G189: a uniform quadratic-seven size bound]
\label{AC-G189}
Set $f=21\operatorname{cond}(\delta)$, with
$\operatorname{cond}(1)=1$.  The group in (G189.1) is finite and
\[
 \frac{\log\#\mathcal S(K,\mathscr D)}b
 \leq
 \begin{cases}
 6\log2,&\delta(-1)=(-1)^e,\\
 \log2+5\mathcal B(e,f),&\delta(-1)=(-1)^{e+1}.
 \end{cases}
 \tag{G189.3}
\]
In particular the second, larger bound is valid in both cases.
It holds on every original coefficient ring, including ramified
ones, without an integral character projector.

For the actual host of (FG.3), let $h\geq2$, let $O_2$ be its
unchanged completion in (FC.2), and let $T\geq1$ be an integer.
There is an absolute constant $C$ such that
\[
 \sum_{\delta\bmod2^{h+2}}\ \sum_{1\leq e<T}
 \frac{\log\#\mathcal S_{2,\delta\xi}(e)}{b_2}
 \leq C\,2^hT^2\bigl(h+\log(T+2)+1\bigr).
 \tag{G189.4}
\]
The sum includes every two-power character and retains every
original local condition.  The fixed fields used in the proof
do not depend on $h,e,q,d$ or the original completion.
This is an upper bound, not a strict-zero assertion.
The nonzero classes of AC-G186 are compatible with it.
\end{theorem}

\begin{proof}
\proofstep{1. Three small, unconditional class certificates.}
The only finite class computations used are the following.
Here $S$ always means all primes above two, and
$\operatorname{Cl}_S$ is the ordinary ideal class group modulo
the classes of those primes.
Write $R=\mathbb Q(\sqrt3,\sqrt7)$,
$F_0=\mathbb Q(i,\sqrt{-7})$ and
$F_1=\mathbb Q(\zeta_8,\sqrt{-7})$.
\[
\begin{array}{c|c|c|c}
\text{field}&\text{defining polynomial}&
 \text{certified class information}&\text{local }(e_2,f_2)\\ \hline
R&X^4-20X^2+16&\operatorname{Cl}(R)=0&(2,2)\\
F_0&X^4-3X^2+4&\operatorname{Cl}(F_0)=0&(2,1),(2,1)\\
F_1&X^8+X^4+16&\operatorname{Cl}_S(F_1)=0&(4,1),(4,1)
\end{array}
\tag{G189.5}
\]
The portable input \texttt{certify\_three\_subfields.gp} uses
PARI/GP 2.19.1, fresh \texttt{bnfinit} and
\texttt{bnfcertify(b,1)} on each literal polynomial.
Each certificate returns one.  The computed class groups are
respectively $1,1,C_2$.  In the last case each prime above two
has exponent vector $(1)$ against the computed class generator;
the displayed principal factor is checked by exact ideal-HNF
multiplication.  Thus that generator is killed in the $S$-class
quotient.  The manual's flag-one conclusion is exactly that the
actual class group is a quotient of the computed group.  This
suffices for every assertion in (G189.5).  No complete unit basis
or equality with $C_2$ is asserted.

The field identifications are exact.  If $x=\sqrt3+\sqrt7$, then
the first polynomial vanishes, and
$\sqrt3=(x-4/x)/2$, $\sqrt7=(x+4/x)/2$.
For a root $x$ of the second polynomial, put
$s=2x^2-3$ and $i=x/(2-x^2)$; then $s^2=-7$, $i^2=-1$.
Conversely $x=i(1-s)/2$ generates the indicated biquadratic field.
For a root $x$ of the third polynomial, put
\[
 s=(2x^4+1)/3,\qquad z=3x/(x^4+2).
\]
Then $s^2=-7$ and $z^4=-1$, so $z$ is primitive of order eight.
Conversely $x=z(1+s)/2$ gives the same polynomial.
The indicated fields have degrees four, four and eight by their
distinct quadratic subfields and ramification.  These identities
therefore identify the full fields, not just subfields.

\proofstep{2. Exact towers and their local decomposition.}
For $n\geq0$ put
\[
 \begin{aligned}
 F_n&=\mathbb Q(\mu_{2^{n+2}},\sqrt{-7}),&
 M_n&=K F_n,\\
 R_n&=\mathbb Q(\sqrt3,\sqrt7)\,
                      \mathbb Q(\mu_{2^{n+2}})^+,&
 M_\infty&=\bigcup_n M_n.
 \end{aligned}\tag{G189.6}
\]
Thus $R_n=M_n^+$, $[F_n:\mathbb Q]=2^{n+2}$,
$[M_n:\mathbb Q]=2^{n+3}$, and $[R_n:R]=2^n$.
The odd ramification of the quadratic factors proves the required
global disjointness from the two-power cyclotomic fields.

The identity $-7\in(\mathbb Q_2^\times)^2$, proved by the
elementary binary Hensel lift in G186, shows that $F_n$ has two
places over two, each with completion $\mathbb Q_2(\mu_{2^{n+2}})$.
The field $K_2$ is the unramified quadratic extension of
$\mathbb Q_2$.  Hence both places are inert in $M_n/F_n$,
and $M_n$ again has exactly two places over two.
All these places are totally ramified in the respective
cyclotomic towers.

The field $R$ has one place $\mathfrak p$ over two.  More explicitly,
$\mathbb Q_2(\sqrt7)=\mathbb Q_2(i)$ because $-7$ is a square,
and $\mathbb Q_2(\sqrt{21})$ is the unramified quadratic
extension because $21/5$ is a square in $\mathbb Q_2$.
Thus
\[
 R_{\mathfrak p}=\mathbb Q_2(i)\,\mathbb Q_2^{\rm unr,2}.
\]
The completion of $R_n$ is
$\mathbb Q_2(\mu_{2^{n+2}})\mathbb Q_2^{\rm unr,2}$;
its degree equals $[R_n:\mathbb Q]$.  There is one place
at every layer, totally ramified over $R_{\mathfrak p}$.

The quadratic extension $M_n/F_n$ ramifies only at primes above
three.  Those primes are totally ramified, since $F_n/\mathbb Q$
is unramified at three and $K/\mathbb Q$ is ramified there.
There are at most four such primes.  Indeed Frobenius three in
\[
 \operatorname{Gal}(F_n/\mathbb Q)
   \simeq(\mathbb Z/2^{n+2}\mathbb Z)^\times\times C_2
\]
is $(3,-1)$.  At $n=0$ its order is two, giving two primes.
For $n\geq1$, the exact order of three modulo $2^{n+2}$ is
$2^n$, by the elementary lifting-of-exponent formula, and the
second coordinate has order two; hence there are four primes.
There is no archimedean ramification because $F_n$ is imaginary.

\proofstep{3. The two integral norm relations.}
First, every $R_n$ has trivial two-primary ordinary class group.
Here is a direct one-prime argument, including the real places.
Let $H/R_n$ be its maximal unramified abelian two-extension
with every real place split.  It is Galois over $R$.
If $H\ne R_n$, its inertia subgroup over the unique prime of
$R$ above two has order $[R_n:R]$ and is proper in the finite
two-group $\operatorname{Gal}(H/R)$.
A maximal subgroup containing it is normal of index two.
The corresponding quadratic extension of $R$ is unramified at
every finite place and split at every real place, since $H$ is
totally real.  This contradicts $\operatorname{Cl}(R)=0$.
Thus the asserted class-group vanishing holds for every $n$.

Second, the two-primary $S$-class group of $F_n$ vanishes for
every $n$.  We give the stabilization argument at the exact
scope needed here.  It is the special case of Coates--Li--Li,
\emph{Classical Iwasawa theory and infinite descent on a family
of abelian varieties}, Propositions 3.2--3.3 and Lemma 3.5,
whose complete applied proof is retained in the source bundle.
Fix $N\geq1$ and the finite layer $F_N/F_0$.  Let $H_N/F_N$ be its maximal
unramified abelian two-extension in which all primes above two
split.  Write $Y=\operatorname{Gal}(H_N/F_N)$ and let
$\gamma$ generate $\operatorname{Gal}(F_N/F_0)$.
Each decomposition group at a prime above two maps isomorphically
onto this cyclic group.  Choose generators $\sigma_1,\sigma_2$
lifting the same $\gamma$, and set $y=\sigma_2\sigma_1^{-1}\in Y$.
The commutator subgroup is $(\gamma-1)Y$.
Killing the decomposition groups and commutators identifies the
$S$-class group at level zero with $Y/Z$, where
\[
 Z=(\gamma-1)Y+\mathbb Z_2[\gamma]y.
\]
At level one the same argument gives $Y/(1+\gamma)Z$:
indeed $\sigma_2^2\sigma_1^{-2}=(1+\gamma)y$ and
$\gamma^2-1=(1+\gamma)(\gamma-1)$.
Both quotients vanish by (G189.5).  Therefore
$Y=Z=(1+\gamma)Z$.
The element $1+\gamma$ belongs to the maximal ideal
$(2,\gamma-1)$ of the finite two-group ring
$\mathbb Z_2[\operatorname{Gal}(F_N/F_0)]$.
Nakayama's lemma gives $Y=0$.
This proves the assertion for every $N$; the zero layer is already
in (G189.5).  All decomposition hypotheses used here were checked
in Step 2, and no assertion about an unrelated split-prime tower
has been substituted.

Put $A_n=\operatorname{Cl}_S(M_n)\{2\}$.
Let $c$ be complex conjugation and let $r$ fix $F_n$ and send
$\sqrt{-3}$ to its negative.  Extension after norm on $S$-class
groups is the sum of the two automorphisms.  The two vanishings
just proved imply the integral identities
\[
                  (1+c)A_n=0,\qquad(1+r)A_n=0.
 \tag{G189.7}
\]
No division by two has occurred.

\proofstep{4. At most five generators in the fixed infinite tower.}
We bound $\dim_{\mathbb F_2}A_n/2A_n$ uniformly, using the
elementary proof of Proposition 3.4 of the same primary source.
For the quadratic extension $M_n/F_n$, take the two primes in
$S$ and all primes ramified above three.  Their number $t$ is
at most six.  In the Galois group over $F_n$ of the $S$-Hilbert
two-class field of $M_n$, use decomposition groups at the two
$S$-primes and inertia groups at the ramified primes.
Each maps isomorphically onto $\operatorname{Gal}(M_n/F_n)$.
Choose lifts $\sigma_1,\ldots,\sigma_t$ of its nontrivial element
and put $y_i=\sigma_i\sigma_1^{-1}\in A_n$ for $i>1$.
As above the commutator subgroup is $(r-1)A_n$.
The maximal abelian quotient that kills all these local groups
is the $S$-class two-group of $F_n$, which is zero.
Consequently
\[
 A_n=(r-1)A_n+
             \sum_{i=2}^t\mathbb Z_2[r]y_i.
\]
By (G189.7), $r=-1$ on $A_n$, so the images of the $t-1$
elements $y_i$ generate $A_n/2A_n$.  Hence
\[
                   \dim_{\mathbb F_2}A_n/2A_n\leq5.
 \tag{G189.8}
\]

The cyclotomic norm maps $A_{n+1}\to A_n$ are surjective.
The $S$-Hilbert two-class field of $M_n$ is split at its two
$S$-primes, whereas $M_{n+1}/M_n$ is totally ramified there.
Their intersection is therefore $M_n$.
Base change of that class field is an unramified, $S$-split
abelian two-extension of $M_{n+1}$; restriction of Artin maps
is precisely the norm on the class groups, proving surjectivity.
Set
\[
                      X=\varprojlim_n A_n.
 \tag{G189.9}
\]
The quotient $X/2X$ is the inverse limit of $A_n/2A_n$.
One can check exactness directly: compatible elements in $2A_n$
have compatible halves by compactness of the finite nonempty
sets of choices; every finite set of compatibility requirements
is solved by taking a half at its highest level and norming down.
The same finite-choice argument lifts compatible residues modulo
two.  Thus (G189.8) implies $\dim_{\mathbb F_2}X/2X\leq5$.
Lift a basis and take its closed $\mathbb Z_2$-span $Y$ in $X$.
Then $(X/Y)/2(X/Y)=0$.  A nonzero compact pro-two
$\mathbb Z_2$-module has a nonzero finite quotient, whose quotient
modulo two is nonzero.  Therefore $X/Y=0$.
We have proved that $X$ has at most five $\mathbb Z_2$-generators.
This conclusion uses neither a qualitative $\mu=0$ assertion
nor a bound on arbitrary specialization roots.
The two relations (G189.7) pass to $X$.

\proofstep{5. Strict restriction, with the finite descent cost retained.}
The action on $\mathscr D$ is trivial over $M_\infty$.
Restriction of a strict cocycle there is consequently a continuous
homomorphism, unramified at every finite place and zero on the
whole decomposition groups above two.
It therefore factors through $X$ of (G189.9).
For clarity, the infinite-level assertion can be verified directly.
The finite image of such a homomorphism cuts out a finite abelian
two-extension of $M_\infty$ stable under its cyclotomic Galois group.
After increasing the finite layer, conjugation on this finite
kernel is trivial.  The resulting central extension of
$\mathbb Z_2$ by a finite group splits by lifting a topological
generator: its closed procyclic subgroup maps isomorphically
onto $\mathbb Z_2$.
It thus descends to a finite abelian extension at that layer.
Away two it is already unramified, since the cyclotomic tower is
unramified there.  At each of the two places above two its
decomposition subgroup is a graph from $\mathbb Z_2$ into the
finite kernel.  Both graphs vanish on a sufficiently small open
subgroup.  Increasing the layer again makes the descended
extension split at both places.  It belongs to the finite
$S$-class field, giving exactly the factorization through (G189.9).

Write
\[
 G=\operatorname{Gal}(M_\infty/K)
       =U\times\langle\tau\rangle,
       \quad U=\mathbb Z_2^\times,\quad \xi(\tau)=-1.
 \tag{G189.10}
\]
The two odd quadratic ramification loci give this product literally.
The image of $G_{K_2}$ is exactly $U\times\{1\}$:
$\sqrt{-7}\in\mathbb Q_2$, and the unramified quadratic $K_2$
is disjoint from the totally ramified two-power cyclotomic tower.
Inflation--restriction is injective at its first term, globally
and locally.  Hence the kernel of strict restriction to $X$
is contained in
\[
 \begin{aligned}
 \ker\bigl(H^1(G,\mathscr D)\to H^1(U,\mathscr D)\bigr)
       &=H^1(\langle\tau\rangle,\mathscr D^U)\\
       &=\mathscr D^U/2\mathscr D^U.
 \end{aligned}
 \tag{G189.11}
\]
The last equality uses the action $\tau=-1$, with norm zero.
The invariant module $\mathscr D^U$ is a finite cyclic
$O$-module.  Indeed its annihilator is the nonzero ideal generated
by $\delta(u)u^{-e}-1$: on an open subgroup $\delta=1$, and
$e>0$ makes this scalar nontrivial.  Thus
\[
              \#(\mathscr D^U/2\mathscr D^U)\leq2^b.
 \tag{G189.12}
\]
This is the finite descent cost; it has not been discarded.
The image lies in $\operatorname{Hom}_G(X,\mathscr D)$.

In the full rational Galois group, $c r$ fixes $K$ and is the
element $(-1,\tau)$ of (G189.10).  By (G189.7) it acts as one
on $X$.  Its scalar on $\mathscr D$ is
$\delta(-1)(-1)^{e+1}$.
If $\delta(-1)=(-1)^e$, every equivariant homomorphism is killed
by two.  Evaluation on at most five generators of $X$ embeds
their group into $(E/O)[2]^5$, of order $2^{5b}$.
Together with (G189.12) this proves the first case of (G189.3).

Suppose henceforth that $\delta(-1)=(-1)^{e+1}$.
Because $r=-1$ on $X$, every such homomorphism is automatically
equivariant for the full rational character
\[
        \alpha=\chi c_2^{-e},\qquad
        \chi=\delta\xi\epsilon_K=\delta\chi_{21}.
 \tag{G189.13}
\]
The character $\chi$ is primitive of conductor $f$ and has parity
$\chi(-1)=(-1)^{e+1}$; the class character $\alpha$ is odd.
This extension is forced by the integral norm relation, not by
an integral projector at two.

\proofstep{6. An integral compatible Stickelberger annihilator.}
Let $m_n=21\cdot2^{n+2}$ and $H_n=\mathbb Q(\mu_{m_n})$.
It contains $M_n$ with degree three.
In its rational group algebra define
\[
 \begin{aligned}
 \theta_n'&=\sum_{\substack{1\leq a<m_n\\(a,m_n)=1}}
           \left(\frac a{m_n}-\frac12\right)\sigma_a^{-1},\\
 Z_n&=(5-\sigma_5)\theta_n'
   =\sum_{\substack{1\leq a<m_n\\(a,m_n)=1}}
            \left(\left\lfloor\frac{5a}{m_n}\right\rfloor-2\right)
                       \sigma_a^{-1}.
 \end{aligned}\tag{G189.14}
\]
The last formula proves integrality directly.
In the convention of Sharifi, Definition 4.5.1, the usual
Stickelberger element is $\theta_n'+N_n/2$.
Lemma 4.5.2 and the complete proof of Theorem 4.5.6 show that
$(5-\sigma_5)(\theta_n'+N_n/2)$ annihilates the class group.
The difference is $2N_n$, which also annihilates it because
extension after norm to $\mathbb Q$ is zero on ideal classes.
Thus $Z_n$ annihilates $\operatorname{Cl}(H_n)$.
Extension from $M_n$ to $H_n$ is injective on the two-primary
class group, since norm after extension is multiplication by
three.  The pushdown of $Z_n$ therefore annihilates $A_n$.

The elements $\theta_n'$ are compatible under group-ring
projection: the two representatives $a,a+m_n$ at level $n+1$
have Bernoulli coefficients summing to $a/m_n-1/2$.
Consequently the integral $Z_n$ form a compatible bounded
measure $Z$ and their pushdowns annihilate $X$.
For every equivariant homomorphism in Step 5 its image is
annihilated by the scalar $A=\alpha(Z)\in O$.

We derive this scalar rather than infer it from a characteristic
ideal.  For $n$ large enough that $\operatorname{cond}(\delta)$
divides $2^{n+2}$, put $k=e+1$ and
\[
 S_{n,k}=\sum_{\substack{1\leq a<m_n\\(a,m_n)=1}}
                            \chi^{-1}(a)a^k.
\]
The elementary Bernoulli-polynomial sum formula gives
\[
 \lim_n\frac{S_{n,k}}{m_n}
      =\beta_e B_{k,\chi^{-1}},\qquad
 \beta_e=\begin{cases}1,&2\mid f,\\1+2^e,&f=21.\end{cases}
 \tag{G189.15}
\]
Indeed, without the restriction to odd integers, splitting the
sum into periods of length $f$ and using
$\sum_{j=0}^{q-1}(j+x)^k=
 [B_{k+1}(q+x)-B_{k+1}(x)]/(k+1)$ gives
$B_{k,\chi^{-1}}$ as $q\to0$ two-adically.
When $f$ is odd, subtracting the even terms contributes
$\chi^{-1}(2)2^{k-1}=-2^e$, giving the displayed factor.

Put $q_a=\lfloor5a/m_n\rfloor$ and $r_a=5a-m_nq_a$.
The $r_a$ permute the units modulo $m_n$, and
$\chi(r_a)=\chi(5)\chi(a)$.
Expand $r_a^k$ and compare the weighted sum with $S_{n,k}$.
Terms of binomial degree at least two, after division by $m_n$,
tend to zero two-adically: $e$ is fixed and every remaining sum
has integral terms.  The linear term gives
\[
 \lim_n\sum_a\chi^{-1}(a)a^e q_a
  =\frac{5^{e+1}-\chi(5)}{(e+1)5^e}\,
                             \beta_e B_{e+1,\chi^{-1}}.
\]
The sum $\sum_a\chi^{-1}(a)a^e$ tends to zero by the same
Bernoulli sum formula, before division by $m_n$.
Evaluating the integral measure (G189.14) therefore yields
\[
 A=(5-\chi(5)5^{-e})\,\beta_e\,
                    \frac{B_{e+1,\chi^{-1}}}{e+1}.
 \tag{G189.16}
\]
The use of integral coefficients in (G189.14) makes this a
literal continuous-measure specialization; it does not reduce
a denominator-containing sum modulo its own precision.
Changing a lift of a residue class changes its character value
by a quantity tending uniformly to zero.  The coefficients are
integral, so the nonarchimedean triangle inequality gives the
same bound for the total Riemann-sum error, without multiplying
by the growing number of residue classes.
There are no omitted Euler factors at three or seven, since
both divide the primitive conductor $f$.

\proofstep{7. Nonzero algebraic value and its global height bound.}
For completeness, the Bernoulli value in (G189.16) is nonzero
and has a uniform bound at every complex embedding.
For $k=e+1\geq2$, the absolutely convergent Fourier series
\[
 B_k(x)=-\frac{k!}{(2\pi i)^k}
                  \sum_{n\ne0}\frac{e^{2\pi i n x}}{n^k}
\]
follows by integration by parts of the periodic Bernoulli
polynomial; its nonzero Fourier coefficients are
$-k!/(2\pi i n)^k$ and its mean is zero.
Insert this into
$B_{k,\chi^{-1}}=f^{k-1}\sum_{a=1}^f\chi^{-1}(a)B_k(a/f)$.
The primitive Gauss-sum identity and
$\chi(-1)=(-1)^k$ give
\[
 \left|\frac{B_{e+1,\chi^{-1}}}{e+1}\right|
  =\frac{2e!f^{e+1/2}}{(2\pi)^{e+1}}
                      |L(e+1,\chi)|
  \leq4e!f^{e+1/2}.
 \tag{G189.17}
\]
The Gauss sum has absolute value $\sqrt f$, and the Euler
product at $e+1>1$ is nonzero and bounded by $\zeta(e+1)\leq2$.
These facts follow respectively by primitive finite character
orthogonality and the absolutely convergent Euler product.
They prove nonvanishing as well as the bound.
The smoothing factor in (G189.16) is nonzero because
$5^{e+1}$ cannot be a root of unity, and $\beta_e$ is nonzero.
Thus $A\ne0$.

Let $C=\mathbb Q(\delta)$, a two-power cyclotomic field, possibly
$\mathbb Q$.  It has a unique prime above two and contains $A$.
The recurrence for Bernoulli numbers proves by induction that
$(j+1)!B_j$ is integral.  Expanding $B_{k,\chi^{-1}}$ in powers
of $a/f$ then shows
$f(k+1)!B_{k,\chi^{-1}}\in O_C$.
Consequently the positive integer
\[
                    d=(e+1)f(e+2)!5^e
\]
has $dA\in O_C$.
At every complex embedding, (G189.16)--(G189.17),
$|5-\chi(5)5^{-e}|\leq6$ and $|\beta_e|\leq2^{e+1}$ give
\[
 |dA|\leq48\cdot10^e(e+1)(e+2)!e!f^{e+3/2}.
\]
The nonzero integer norm of $dA$ and the unique prime above two
therefore imply, with $v_2(2)=1$,
\[
 v_2(A)\log2\leq v_2(dA)\log2
       \leq\frac{\log|N_{C/\mathbb Q}(dA)|}{[C:\mathbb Q]}
       \leq\mathcal B(e,f).
 \tag{G189.18}
\]
The middle inequality uses divisibility of a nonzero integer
norm by its exact power of two; all embeddings over two have
the same normalized valuation because that prime is unique.
This bounds an actual algebraic annihilator and does not assume
a lower bound for a higher regulator or an unspecified period.

\proofstep{8. Return to the unchanged original lattice and sum the family.}
Since $X$ has at most five $\mathbb Z_2$-generators and $A$
annihilates the image of every equivariant map, evaluation embeds
that group of maps into $((E/O)[A])^5$.
For the original $O$ one has exactly
\[
 \log\#(E/O)[A]=[E:\mathbb Q_2]v_2(A)\log2.
\]
This identity includes the ramification index of $O$; it does
not replace normalized length by residue degree alone.
Combining (G189.12), (G189.18) and the restriction map proves
the second case of (G189.3), including finiteness.
Since $\mathcal B(e,f)\geq\log2$, it also bounds the first case.

In the actual host every $\delta\xi$ has values of order dividing
$2^h$, already in its original $O_2$.
Its coefficient is literally
\[
 (O_2\otimes_{\mathbb Z_2}\mathbb Q_2/\mathbb Z_2)
                         (\delta\xi)(-e)
       =(E_2/O_2)(\delta\xi c_2^{-e}).
 \tag{G189.19}
\]
Inflation from the actual ramification group of (FC.2) is
injective on this module.  Its full local condition above two
and every away inertia condition give the strict class used
above; added host primes remain present.
Thus (G189.3) bounds each actual group on its original ring.

There are $2^{h+1}$ characters modulo $2^{h+2}$, and their
primitive conductors are at most $2^{h+2}$.
The factorial bound $j!\leq(j+1)^j$ gives an absolute $C_1$ with
\[
 \mathcal B(e,21\cdot2^{h+2})
       \leq C_1(e+2)\bigl(h+\log(e+2)+1\bigr).
\]
Sum (G189.3) over $1\leq e<T$ and all these characters.
For $T=1$ the sum is empty; for $T\geq2$, the sums of $e+2$
and of $(e+2)\log(e+2)$ have the stated quadratic bounds.
This proves (G189.4).
Since $D_u=144\cdot2^h$ in (FC.1), its right side is bounded
by an absolute multiple of $D_u^2T^2\log(8TD_u)$.
On $T=\lfloor u^{399}\rfloor$ it has the same polynomial
envelope permitted by (FL.4).  A separate union proof must still
add it to that one full local-error majorant and retain the same
window before changing an omitted set.
\end{proof}

\paragraph{Source and scope boundary.}
The exact class certificates and ideal identities, including their
limited flag-one meaning, are preserved with the portable input,
outputs and original PARI manual passages.  The failed larger-field
stabilization exploration is preserved separately and is not an input.
The complete applied group proofs of Coates--Li--Li,
Propositions 3.2--3.4 and Lemma 3.5, are read and specialized
above; the source is Selecta Mathematica 27 (2021), article 28,
\url{https://doi.org/10.1007/s00029-021-00644-3}.
The complete applied Sharifi Stickelberger proof, with the
ordinary arithmetic Frobenius convention, is retained at
\url{https://www.math.ucla.edu/~sharifi/notes/iwasawa-ch04.html},
Definition 4.5.1, Lemmas 4.5.2--4.5.5 and Theorem 4.5.6.
Hilbert class-field reciprocity and elementary inflation--restriction
retain the previously audited source scopes; no new proof of the
foundations of class field theory is claimed.
The moment and global norm bounds are derived above.
No main conjecture, uniform conjectural regulator estimate,
integral division by two, coefficient replacement, arbitrary
near-one $\kappa$ extension, or Lean execution is used.
The actual finite-cover/Artin--Tate interpretation is the retained
SGA 1, Expos\'e V, sections 4 and 7, and Stacks 03PK dictionary.
Only the pure quadratic-seven coefficient-two family is bounded;
the mixed and auxiliary families, full numerical remainder,
$\mathcal R_{\rm cub}$ and H04.4 remain separate.
